\textbf{COLUNI 2025 -- Questão 07}\par\medskip Leia o trecho do poema \textit{Meu Rosário}.

E sonho nas contas de meu rosário lugares, pessoas,\\
vidas que pouco a pouco descubro reais.\\
Vou e \underline{volto} por entre as contas de meu rosário,\\
que são pedras marcando-me o corpo-caminho.\\
E neste \underline{andar} de contas-pedras,\\
o meu rosário se \underline{transmuda} em tinta,\\
me guia \underline{o dedo},\\
me insinua \underline{a poesia}.

Com base no trecho e nas palavras destacadas, analise as afirmativas a seguir.

I. \underline{andar}: verbo flexionado no presente do indicativo.\\
II. \underline{o dedo}: expressão que funciona como complemento nominal.\\
III. \underline{volto}: verbo flexionado na primeira pessoa do singular.\\
IV. \underline{a poesia}: expressão que funciona como objeto direto.

Os termos destacados estão CORRETAMENTE classificados em:
\begin{enumerate}
\item I e IV.
\item I, apenas.
\item III e IV.
\item II, apenas.
\end{enumerate}
